\documentclass[11pt]{article}
\usepackage[letterpaper,margin=0.85in]{geometry}
\usepackage[T1]{fontenc}
\usepackage{lmodern,microtype,amsmath,amssymb,booktabs,tabularx,array,longtable}
\usepackage{xcolor,enumitem,fancyhdr}
\usepackage[hyphens]{url}
\usepackage{hyperref}
\definecolor{navy}{HTML}{17324D}
\definecolor{teal}{HTML}{176B73}
\hypersetup{colorlinks=true,linkcolor=navy,citecolor=teal,urlcolor=teal,
 pdftitle={Interactive Batch ORAM with Garbled Circuits and Secret-Shared Shuffling},
 pdfauthor={Technical design and cost study}}
\setlength{\parindent}{0pt}
\setlength{\parskip}{5pt}
\setlength{\emergencystretch}{2em}
\setlength{\headheight}{14pt}
\setlist{nosep,leftmargin=1.6em}
\pagestyle{fancy}\fancyhf{}
\fancyhead[L]{\small\textcolor{navy}{Interactive batch ORAM}}
\fancyhead[R]{\small Scheme and cost derivation}
\fancyfoot[C]{\thepage}
\newcommand{\AND}{\mathbin{\wedge}}
\newcommand{\xor}{\mathbin{\oplus}}
\newcommand{\ceil}[1]{\left\lceil #1\right\rceil}
\newcommand{\bits}{\operatorname{bits}}
\newcommand{\valid}{\operatorname{valid}}
\newcommand{\uid}{\operatorname{uid}}
\newcommand{\bytes}{\operatorname{bytes}}
\newcommand{\costtable}[2]{\begin{center}\small #2\end{center}\par\textit{#1}\par}
\newcommand{\PrimaryAnnualSmall}{\$39.05M}
\newcommand{\PrimaryAnnualLarge}{\$51.40M}

\begin{document}
{\color{teal}\large TECHNICAL REPORT\par}
{\color{navy}\LARGE\bfseries Interactive Batch ORAM with Garbled Circuits and Secret-Shared Shuffling\par}
\vspace{4pt}
\textbf{Scheme, security conditions, and reproducible resource estimates}\par
16 October 2026 \quad $\cdot$ \quad Two-party semi-honest setting
\vspace{8pt}

\begin{abstract}
We study a private key--value service combining two-choice cuckoo hashing,
recursive Circuit-OPRAM-style batching, streamed garbled circuits, and an
interactive secret-shared shuffle. Both the database and each 256-bit query
key are private; neither server is assumed to know the query. An AES-CMAC
PRF computes both cuckoo indices inside the secure computation. Our main
optimization matches requests to pool records, attaches secret request tags,
shuffles the tagged records with a permutation unknown to either party, and
opens only the shuffled tags. Payload extraction and removal then become
local operations on shares. A fixed number of uniquely tagged dummies hides
pool misses and duplicate requests.

We derive a complete accounting ledger for the selected compilation template,
including garbled tables, both directions of OT traffic, shuffle preprocessing,
conversion between representations, recursive-map updates, persistent state,
and modeled storage I/O. The primary estimate uses ordinary Boolean records
outside the shuffle; it does not require an unverified two-label conditional
swap. Packed-value storage is reported as a conditional sensitivity.
At 10,000 lookups/s, the selected batch has 32 requests for both studied sizes.
The resulting historical-price network, state-capacity, and modeled-I/O budgets
are \PrimaryAnnualSmall{} and \PrimaryAnnualLarge{} per year. These are
capacity scenarios, not measured deployment costs or throughput guarantees:
the inherited parallel eviction and grouping routines remain circuit
composition allowances, and CPU, state service time, initialization, and
overflow probabilities require validation.
\end{abstract}

\begin{center}\small
\begin{tabular}{lrr}\toprule
Primary selected scenario & $N=2^{28}$ & $N=2^{32}$\\\midrule
Batch / map fanout & 32 / 4 & 32 / 4\\
Ordinary ANDs/query (million) & 6.997 & 8.867\\
Total framed traffic/query (MB) & 237.832 & 301.293\\
NIC hosts per party & 298 & 378\\
Persistent state per party (TB) & 64.70 & 1,061.58\\
Network + capacity + modeled I/O/year & \$39.05M & \$51.40M\\
\bottomrule\end{tabular}
\end{center}

\textbf{Interpretation.} The report provides a concrete pool-fetch protocol
and a reproducible cost model for the surrounding construction. It does not
claim that the whole construction has been implemented or that the chosen
capacity parameters have a service-lifetime failure bound.

\clearpage
{\small\setlength{\parskip}{0pt}\tableofcontents}
\clearpage

\section{Functionality, threat model, and notation}
\label{sec:scope}
Two noncolluding servers, a garbler $G$ and evaluator $E$, hold a private
key--value dictionary. We assume a static, semi-honest adversary corrupting
at most one server. A query key $x\in\{0,1\}^{256}$ is supplied as
$x=x_G\xor x_E$; neither share alone reveals $x$. Database initialization
must similarly establish a secret shared or encoded dictionary. Both
cuckoo probes execute on every lookup. The returned value and hit bit remain
shared between the servers; an authorized client may reconstruct them through
a separately specified delivery interface. No server receives the cleartext
query, cuckoo index, or value by default.

This report models application \emph{reads}. Each read still changes the ORAM
position maps, locations, validity bits, and stash/pool state. Arbitrary
application writes, concurrent database mutation, client authentication,
malicious security, availability, and durable recovery need additional
protocols. Public workload information includes the fixed batch size,
provisioned capacities, execution schedule, and completed-query count.

\begin{center}\small
\begin{tabularx}{\linewidth}{lX}\toprule
Symbol & Meaning\\\midrule
$N=2^d$ & Number of dictionary records; $d\in\{28,32\}$\\
$Q$ & Aggregate rate, $10{,}000$ complete dictionary lookups/s\\
$B$ & Fixed batch size; both tables receive $B$ requests\\
$h_\ell$ & Logical-address/leaf width at recursive tree $\ell$\\
$F=2^f$ & Number of child position labels packed in one map record\\
$Z$ & Tree bucket capacity, fixed at 3 in this study\\
$\sigma$ & Additional pool slack, fixed at 128 records\\
$r=\log_2(2B)$ & Number of upper tree levels represented by the pool\\
$P=Z(2^r-1)+\sigma$ & Persistent pool capacity per recursive tree\\
$D_\ell$, $W_\ell$ & Payload width and complete record width in bits\\
$n=P+B$ & Number of tagged records shuffled per tree and batch\\
$\lambda$ & Garbling security parameter, 128 bits\\
\bottomrule\end{tabularx}
\end{center}
All MB, GB, TB, and PB in traffic/state tables are decimal. GiB and MiB are
binary units and are used where required by storage billing. Costs include
both table probes unless a row is explicitly labeled ``per table.''

\section{Dictionary and recursive ORAM layout}
\label{sec:layout}
We use two cuckoo tables, each with $2N$ logical cells, together with an
eight-entry application cuckoo stash. The total table load is $N/(4N)=1/4$.
The original dictionary occupies 16 GiB or 256 GiB before hash-table, ORAM,
and encoding overhead. A data payload has 256 key bits, 256 value bits, and
one occupancy bit:
\begin{equation}
 D_0=513,\qquad h_0=d+1,\qquad W_\ell=D_\ell+2h_\ell+1.
\end{equation}
The two $h_\ell$ fields are the logical record identifier and assigned leaf;
the last bit is the ORAM validity flag. Dictionary occupancy and ORAM
validity are distinct. Recursive position-map records hold $F$ leaf labels:
\begin{equation}
 h_{\ell+1}=h_\ell-f,\qquad D_{\ell+1}=Fh_\ell.
\end{equation}
Map packing changes record width; it does not change the binary tree's
branching factor.

At each tree, the top $r$ levels are represented by an unordered pool.
Below them, a path contains
\begin{equation}
 L_\ell=h_\ell-r+1
\end{equation}
bucket levels. Public old-leaf paths determine which lower buckets are
fetched. An ORAM logical identifier, an assigned leaf, and a physical slot
are different objects. Revealing a requested old leaf does not reveal the
record's identifier or its position within a bucket or pool.

Recursion stops at a configurable base cutoff or before the next tree would
have fewer than $2B$ leaves. The final scalar position map has
$C=2^{h_{\rm last}}$ entries of width $h_{\rm last}$; its \emph{full} capacity,
not the requested cutoff, is charged. The selected layouts and all recursive
record widths are listed in Section~\ref{sec:results}.

The overall structure follows Circuit OPRAM~\cite{chan2017}. The concrete
Boolean compilation below is a source-inspired planning template. In
particular, $Z=3$ and $\sigma=128$ are fixed study parameters, not a proven
overflow bound for these database sizes and access horizons.

\section{Cryptographic primitives and representation costs}
\label{sec:shuffle-primitives}

\subsection{Hashing a secret-shared 32-byte query key}
The query key is shared between the garbler and evaluator:
$x=x_G\oplus x_E\in\{0,1\}^{256}$.  Neither party is assumed to know $x$.
Consequently, computing its two table addresses locally is not an admissible
optimization.  A concrete baseline evaluates AES-128-CMAC inside the secure
computation, with a fixed, separately sampled secret PRF key $K$ shared between
the parties.  The database uses the same function when it is constructed.
For the fixed-length input $x=M_1\mathbin\Vert M_2$, CMAC is
\begin{align}
 L&=\operatorname{AES}_K(0^{128}), &
 K_1&=\operatorname{dbl}(L),\nonumber\\
 C_1&=\operatorname{AES}_K(M_1), &
 T&=\operatorname{AES}_K(C_1\oplus M_2\oplus K_1).
 \label{eq:shared-cmac}
\end{align}
Both message blocks are complete, so CMAC uses $K_1$, without a padding block.
Multiplication by $X$ in the binary field, denoted $\operatorname{dbl}$, consists
of a shift and XORs with selected bits of the constant reduction polynomial;
it requires no AND gates.  These choices follow AES and the CMAC specification
\cite{aesfips,cmac}.

For an auditable gate count, use the established 32-AND AES S-box circuit listed
by NIST, rather than assuming a newer optimized implementation.  The ten AES
rounds contain $160$ S-boxes, giving $160\cdot32=5120$ ANDs for the encryption
datapath.  The AES-128 key schedule contains another $40$ S-boxes, or $1280$
ANDs.  NIST lists $6400$ ANDs for AES-128 with its key schedule and $5120$ for
the fixed-key datapath \cite{nistcircuits}.  A secret, already expanded key has
the same nonlinear datapath count: AddRoundKey remains XOR.  Reusing the
expanded key and $K_1$ across the $B$ queries of a single batch therefore gives
\begin{equation}
 A_{\rm hash}(B)=
 \underbrace{1280+5120}_{6400\text{ setup}}
 +\underbrace{2\cdot5120B}_{10240B\text{ query work}}.
 \label{eq:shared-hash-cost}
\end{equation}
At 32 bytes per half-gate AND, this is $200$ KiB of setup plus $320$ KiB per
query.  The primary model retains the 11 expanded round keys and $K_1$ as
1536 ordinary encoded state bits, using the same linked-garbling interface
as other persistent state.  Setup is therefore performed once per initialized
state and amortized, rather than repeated in every batch.  This requires
compatible labels and distinct gate/hash tweaks across the composed
computation; it does not permit reuse of one-time garbled tables.  A reset to
an incompatible encoding must pay for secure key-state translation or new
input and setup.  Equation~\eqref{eq:shared-hash-cost} counts neither initial
database hashing nor interface traffic.

One 128-bit PRF result suffices for two $h$-bit table addresses when $2h\le128$:
take two disjoint $h$-bit slices of $T$.  For a pseudorandom $T$, the slices are
jointly pseudorandom; a second CMAC is unnecessary for this address-generation
interface.  Equal keys deliberately receive equal address pairs.  The full
query key must still be compared with candidate keys to handle address
collisions.  This construction assumes fixed-length, domain-separated use of
the PRF; changing its application domain must preserve that separation.

Query-share input and result-share output belong in a separate interface ledger.
Reconstructing $x_G\oplus x_E$ inside GC costs XORs but not ANDs.  The model
charges 4096 bytes for G's share labels and 12,288 bytes for E's input OT,
giving 16,384 bytes per query: 12,288 forward and 4096 reverse, before framing.
This is a convention, not an unavoidable lower bound: existing garbled inputs
and correlated-OT encodings can change the charge.  The value and hit flag
remain shared; taking the two parties' label color bits gives local XOR
shares as explained in Section~\ref{sec:interactive-pool-protocol}.
There is no plaintext answer decoding to E.  Client delivery is a separate
boundary and is not included in the inter-server totals.

\subsection{Packed storage and conditional movement are separate costs}
The proposed packed representation carries a 32-byte value in two 128-bit
objects.  We retain this as a payload-storage and transport optimization.
It does not, by itself, specify the protocol for moving those objects under
a secret control bit.  In ordinary Boolean GC, a 128-bit label represents
one Boolean wire; storing a payload in two opaque carriers is a different
representation with its own input, output, and movement interfaces.

For this reason the primary total uses a fully specified Boolean movement
cost: swapping a 256-bit payload under one secret bit costs 256 ANDs, with
metadata added at its actual Boolean width.  A conditional packed-backend
column may instead use a measured or proved parameter $c_2$, the communication
for swapping the two packed payload carriers under a secret control.  Storage
size alone does not establish $c_2$, and the conditional column must also
include conversions and relabeling.  When a shuffle-and-reveal step has made
the physical selection public, both parties can select their own shares
locally; that public routing genuinely avoids secret-control payload gates.
This distinction makes the Boolean estimate auditable while keeping the
packed optimization visible as a concrete implementation target.

\subsection{One-hot garbling: what is obtained and what is charged}
Let $P=2^k$ be the padded size of a small array.  Heath's binary-to-one-hot
primitive creates $H(y)$ using $(k-1)\lambda$ bits of garbled material, or
$16(k-1)$ bytes at $\lambda=128$, when the evaluator knows $y$.
The complete framework is not restricted to public true inputs: it masks
private inputs and cancels the masks within secure modules
\cite{heath2021,heath2024}.

In particular, Section~5.1 of \cite{heath2024} encodes a private short integer
$x\in\mathbb Z_{2^k}$ using a mask $\alpha$ and a one-hot encoding of
$y=x+\alpha\bmod 2^k$.  G knows $\alpha$ and E learns $y$, while neither learns
$x$.  For binary input, the stated construction adds $\alpha$ with a
$(k-1)$-AND ripple-carry adder and then invokes binary-to-one-hot.  With the
paper's two-ciphertext Boolean AND implementation, the conversion material is
\begin{equation}
 C_{\rm bin\to masked\ hot}
 =32(k-1)+16(k-1)=48(k-1)\ \text{bytes}.
 \label{eq:masked-hot}
\end{equation}
This excludes input-label transfer, the small reveal/interface bits, and any
additional conversion required by the consuming operation.  Binary and
arithmetic one-hot outputs must not silently be identified: an arithmetic
consumer may require the corresponding upcast or word conversion.  The
expanded vector still has $P$ entries, so local work and working storage grow
linearly in $P$ for the binary construction, with further width factors for
arithmetic words or payload outer products.

Equation~\eqref{eq:masked-hot} is a cost for a \emph{masked-position}
representation, not a canonical $H(x)$ at the same interface.  A lookup into
a mutable secret array must also align that representation with the array
and account for any masking correction.  The free truth-table evaluation
$T(f)^\mathsf{T}H(y)$ uses a public constant table; it is not automatically a
lookup into independently encoded secret array entries.  Similarly, the cheap
one-hot outer product distributes one scalar or vector to one selected slot,
not a different unknown payload at every slot.  These observations do not
rule out a useful module, especially with interaction and shuffled arrays;
they identify the missing interface that must be specified.  The primary
totals therefore claim no generic one-hot discount for dynamic pool accesses.

\subsection{One-AND swaps and a fused flat-array exchange}
For a control bit $s$ and two $w$-bit words $a,b$, define
\begin{equation}
 d=s\mathbin{\wedge}(a\oplus b),\qquad
 a'=a\oplus d,\qquad b'=b\oplus d.
 \label{eq:one-and-swap}
\end{equation}
The AND is bitwise with the same control on all bits.  Thus a swap costs
$w$ ANDs and $3w$ XORs, rather than two independently implemented muxes.
For $P$ flat entries and a valid unique target index, start a carry word with
the new value and conditionally swap it with each array entry at its one-hot
selection bit.  The final carry is the old value and the array contains the
new one.  A straightforward shared-prefix binary decoder uses $P-2$ ANDs
for $P\ge2$, giving the specified circuit cost
\begin{equation}
 A_{\rm flat\ exchange}(P,w)=Pw+P-2,
 \qquad X_{\rm flat\ exchange}(P,w)=3Pw+P-2.
 \label{eq:flat-fused-exchange}
\end{equation}
These are counts for this explicit decoder, not a claim of minimum circuit
size.  A one-entry array needs no selector circuit.  An additional secret
enable, out-of-range handling, or metadata predicates must be charged if
present.  The fused exchange assumes the replacement word is already
available; an arbitrary function of the returned old word needs its own
evaluation and suitable writeback.  For batched repeated indices, a
single unordered aggregate of selectors is insufficient: preserve sequential
semantics or securely group duplicates and define their read/write order.

\section{Interactive pool retrieval: protocol and security conditions}
\label{sec:interactive-pool-protocol}

This section specifies a replacement for \emph{pool retrieval}, conditional on
the surrounding Circuit OPRAM modules maintaining their existing invariants.
It is not a new proof of the complete ORAM construction.  The replacement uses
two-party secret-shared shuffling~\cite{cgp2020,osn2021}: the parties obtain
fresh XOR shares of a common random permutation of the records, while neither
party learns the composed permutation.  Both parties remain semi-honest, and
the application key and intermediate records remain secret shared.  The
ORAM architecture, recursive request processing, and eviction obligations
remain those of the surrounding construction~\cite{chan2017}.

\subsection{Interface and invariants}

Fix one recursive level and one batch.  Let $P$ be the public pool capacity,
$B$ the public batch size, and $u$ the number of UID bits at this level.  A
pool record contains a validity bit, UID, current leaf label, and payload.
The pool is a \emph{multiset stored in an array}; physical positions carry no
meaning needed by subsequent pool processing.  The following are required.

\begin{enumerate}
\item At most one valid record in the union of the pool and the lower tree
has a given UID.  Invalid records may contain arbitrary stale fields.  When
overlapping public paths are fetched, repeated appearances of one physical
tree slot are aliases, not additional independent live records.
\item Existing private conflict processing supplies a representative bit
$r_j$ for each request $j\in\{1,\ldots,B\}$.  There is at most one
representative per UID, and representatives have distinct UIDs.  Duplicate
group membership and result propagation remain secret.
\item Each requested live block is obtained from exactly one of the pool
and its fetched lower path.  Invalid and nonrepresentative requests are
masked.  The surrounding algorithm privately combines these possibilities
and clears each actual source once.
\item Later eviction and cleanup consume records through their secret
metadata, rather than interpreting a pool position as a fixed tree node.
If an implementation instead assigns persistent semantics to pool slots,
it must first change that interface or pay for inverse routing.
\end{enumerate}

These are interface conditions to verify in an implementation, not additional
leakage permissions.  In particular, the parties must not disclose whether
a request hit the pool, the lower tree, or a dummy.

\subsection{Match, tag, shuffle, and extract}

The tag alphabet is $\{0,1,\ldots,B\}$, with zero denoting an unselected row.
A tag has $t=\lceil\log_2(B+1)\rceil$ bits.  Tags identify public request
\emph{slots} in the current processing order; they do not disclose UIDs or
secret duplicate groups.  The following procedure is executed separately
at each recursive level, with fresh shuffling randomness.

\begin{quote}\small
\textbf{Input:} $P$ secret records $x_i$, and $B$ secret requests
$(q_j,r_j)$ in the surrounding circuit's current order.

\textbf{1. Private matching.} For every $i,j$, compute
\[
 e_{ij}=r_j\wedge\operatorname{valid}(x_i)
              \wedge[\operatorname{UID}(x_i)=q_j],\qquad
 h_j=\bigoplus_{i=1}^{P}e_{ij}.
\]
Set the tag of original row $i$ to the bitwise XOR of the public binary
constants $j$ for which $e_{ij}=1$.  There is at most one such $j$.

\textbf{2. Dummy completion.} Append one invalid record $d_j$ per request.
Give $d_j$ tag $j$ if $h_j=0$, and tag zero otherwise.  Its other fields
have the same fixed width and sharing representation as an ordinary record.
Nonrepresentatives have $h_j=0$ and therefore receive their own dummy.

\textbf{3. Secret shuffle.} Jointly shuffle the $n=P+B$ complete
$(\mathrm{tag},\mathrm{record})$ rows.  Each party samples a private uniform
permutation; the protocol applies their composition and refreshes shares.
Use a fresh protocol instance and fresh correlations for this invocation.

\textbf{4. Open tags only.} Reconstruct the shuffled tag column.  For each
$j$, both parties locally copy their record share at the unique row tagged
$j$ to output slot $j$.  Neither party reconstructs the record or its
validity bit.  Output shares remain secret.

\textbf{5. Retain the complement.} Both parties remove all $B$ nonzero-tag
rows locally and discard the public tags.  Keep the remaining $P$ records
in their current unordered arrangement.  Continue the existing private
path/pool result selection, duplicate-result routing, recursive updates,
and reinsertion/eviction schedule.
\end{quote}

At every invocation, each nonzero tag occurs exactly once: either its real
pool record or its dedicated dummy receives it.  Consequently the number
and pattern of opened tags do not reveal the number of hits.  A single
common dummy would be incorrect, as repeated opened positions would reveal
misses or request equality.  Tag assignment happens \emph{before} shuffling,
so no post-shuffle scan to rediscover dummy positions is needed.

\paragraph{Occupancy and clearing.}
Let $H$ be the number of representative pool hits.  Removal deletes $H$
original valid records and $B-H$ appended dummies.  The remaining $H$
appended dummies replace the removed real records, leaving exactly $P$
rows.  This extracts and clears the pool sources in one operation.  It is
not necessary to privately compact the survivors or return them to their
old slots.  Appending the surrounding algorithm's $B$ refreshed records
and masked dummy returns gives the usual $P+B$ candidate bound.  The
existing additional eviction workspace, candidate selection, and cleanup
are still required; this argument does not reduce their allowances.

\paragraph{Matching cost.}
Equality of two $u$-bit UIDs can be computed with $u-1$ ANDs, followed by
two ANDs for validity and representative gating.  Thus $BP(u+1)$ ANDs suffice
for this matching kernel.  The numerical ledger reserves one additional
AND per pair, charging $BP(u+2)$ as a conservative interface allowance.
Unique-hit aggregation is XOR.  Multiplication
of $e_{ij}$ by a public tag bit is either copying or zero, so forming tags
adds no payload-wide conditional selection.  This count excludes the
retained request-conflict module, path matching, and subsequent result
selection.  If valid-UID uniqueness is not established, XOR is not a
substitute for a general OR or a first-match selection circuit.

\subsection{Composition with recursion, paths, and updates}

Private request sorting remains part of the protocol.  The opened tag
returns a result to the same \emph{processing slot} that generated that tag.
It does not reveal the private sorting permutation.  Returning results to
the original query order still requires the existing private inverse
routing or an independently proved alternative.  Public tag extraction
cannot justify deleting those costs.

At a recursive level, several child requests may refer to one parent
position-map block.  Deduplication therefore uses that level's block UID,
not merely the original application key.  The parent block is fetched once;
its different child offsets are read and updated privately, updates are
coalesced, and one final parent block is committed.  This shuffle does not
make concurrent whole-block writes safe.  For a read batch, duplicate
queries can receive copies of the common result through existing secret
routing.  Supporting application writes additionally requires the chosen
operation order to be implemented privately within each duplicate group.

The pool and lower-tree results must be combined using secret hit/validity
flags.  A no-pool-hit dummy is a placeholder, not evidence of an absent
dictionary key.  Likewise, selecting a real pool row does not authorize
revealing that the lower path was unnecessary.  The existing fixed path
schedule is retained.  Nonrepresentatives use the surrounding scheme's
dummy-path rule, and collisions among public paths do not imply equality
of logical UIDs.  Shared physical slots on overlapping paths must have
their extraction/clear masks aggregated before a single state update.

Fresh leaf labels and updated metadata stay secret.  A removed valid pool
record is reintroduced only through the existing refreshed-return module;
an unselected record survives exactly once.  No later step may both retain
an old live copy and append a fresh live copy of the same UID.  These
conditions, together with correct coalescing of recursive updates, preserve
the valid-UID uniqueness needed by the next invocation.

\subsection{Security argument and its scope}

Consider a hybrid in which matching, conflict processing, share conversion,
and shuffling are replaced by their specified secure functionalities.
Conditional on the prior transcript, the pre-shuffle tag vector contains
exactly one copy of each $j\in\{1,\ldots,B\}$ and $P$ zeros.  A fresh
uniform permutation maps it to a uniform ordered placement of the $B$
distinct nonzero tags.  This distribution is independent of the hit vector,
the secret UIDs, and the original locations.  A simulator can therefore
sample the opened tag vector without learning any of those secrets.

One party knows its own private permutation, but the other party's
independent uniform permutation still hides the composition.  Secure
shuffling must refresh the shares: simply permuting recognizable active
labels or unchanged share strings is not this functionality.  Its output
record shares reveal no values when tags are opened.  Extraction and
complement removal are deterministic local functions of the public tags
and private shares, so they introduce no further openings.  Subsequent
secure modules consume these shares and reveal only outputs already
authorized by the surrounding protocol.  This is a conditional composition
argument for the replacement module, not a security or overflow proof for
the entire Circuit OPRAM template.

Every later batch must use a fresh hidden shuffle before disclosing further
positions.  Reusing a permutation permits cross-batch location linkage;
reusing its masks or correlations can expose differences of secret arrays.
Preprocessing can be performed ahead of time, but each correlation is
consumed once and must not be restored after a retry or rollback.  A retry
uses a new correlation.  In the semi-honest model, ``discard'' means that
the protocol does not reuse a correlation; privacy must hold even against
an adversary retaining its complete past view.

\subsection{Representations and explicit communication charges}

For the ordinary Boolean free-XOR backend, a bit $b$ is represented by a
garbler-held zero label $K^0$ and an evaluator-held active label
$K^b=K^0\oplus b\Delta$, where the permutation/color bit of $\Delta$ is
one.  Taking the color bits locally gives
\[
 [b]_G=\operatorname{col}(K^0),\qquad
 [b]_E=\operatorname{col}(K^b),\qquad
 [b]_G\oplus[b]_E=b.
\]
This adds no communication and discloses nothing outside each party's
existing view.  It does not itself promise independent fresh masks across
all output bits; the secure shuffle supplies the required refresh.  No
active label is sent to the garbler and no zero label is sent to the
evaluator during this conversion.

Returning an XOR-shared bit to an ordinary GC can be charged conservatively
as two party inputs followed by a free XOR: a 16-byte garbler input label,
two 16-byte OT message corrections for the evaluator input, and a 16-byte
per-OT matrix allowance for standard 128-bit IKNP extension.  This gives
\emph{64 bytes per logical re-input bit}, plus amortized base-OT setup and
transport padding.  This is an explicit accounting choice rather than an
optimal conversion claim.  Re-input every retained or extracted field that
the next GC actually consumes; the shuffle cost alone is not the complete
retrieval cost.  Persisting raw XOR shares instead of labels changes the
storage backend and requires charging subsequent touched-state re-input.

The color-bit argument applies to one-label-per-Boolean-bit encodings.
It does \emph{not} turn two opaque packed labels into shares of a 256-bit
value.  Accordingly, an ordinary Boolean model plus the explicit conversion
charge is the concrete baseline.  A two-label payload scenario remains
conditional on a specified packed encoding and secure conversion/update
interface.  One-hot and specialized garbling techniques may improve other
small-array modules~\cite{heath2024,heath2021}, but they do not remove the
representation or composition obligations here.

\paragraph{Shuffle cost.}
Let $R$ include all actual XOR-share bytes in a row, including its tag.
An $n$-input generalized Waksman network has
\[
 S(n)=n\lceil\log_2 n\rceil-2^{\lceil\log_2 n\rceil}+1
\]
switches.  In the construction of~\cite[Appendix A.3.1 and Figure 16]{osn2021},
a switch uses a 1-out-of-2 OT with two $2R$-byte candidate payloads.
Random-OT optimization reduces correction traffic from $4R$ to $2R$
bytes per switch.  Wire masks are chosen in advance, so these OTs can be
prepared in parallel; the data-dependent phase sends one masked input
array per private permutation.  Composing two private permutations gives
\[
 C_{\rm shuffle}=\underbrace{2nR}_{\rm online}
             +\underbrace{4S(n)R}_{\rm switch\ corrections}
             +C_{\rm OT}.
\]
Uniform permutations must be sampled first and then routed through the
network; independent fair switch bits do not sample uniform permutations.
All preprocessing communication belongs in the bandwidth/cost ledger even
when it lies outside the batch's critical path.  For $P=317$, $B=32$, and
$n=349$, $S(n)=2630$.  If a complete tagged row occupies 73 bytes, the first
two terms are 50,954 and 767,960 bytes per batch.  An illustrative 16-byte
OT-extension allowance per switch adds 84,160 bytes across the two
networks: 903,074 bytes per batch, or 28.22 decimal KB per request.  If the
73-byte width excludes the tag, rounding to 74 bytes gives 28.57 KB per
request instead.  Opening the tag column to both parties adds at most
$2\lceil nt/8\rceil$ bytes with bit packing: 524 bytes per batch here,
or 698 bytes if each tag is padded separately to one byte.  Metadata
matching, re-input, result routing, and the remaining ORAM work are
separate charges.


\section{Full batch execution and state dependencies}
\label{sec:pipeline}
The two cuckoo tables execute in parallel, while each table's position-map
chain is dependent. The logical schedule is:
\begin{enumerate}
\item Input the two shares of every query key; evaluate the secret-key PRF
and retain both hash indices as encoded values.
\item Privately group duplicate logical requests at each recursive level.
Do not reveal repeated old leaves directly. Representatives use their old
positions; nonrepresentatives use independently random dummy paths.
\item Access and update the flat base map. Its returned old labels enable
the next recursive tree to disclose its duplicate-masked old paths.
\item For that tree, scan lower-path metadata and perform the pool protocol.
Combine the pool and path results under the one-live-copy invariant.
Forward each representative result to duplicate requests privately, then
restore request order using saved secret routing controls.
\item Coalesce updates to different offsets of the same recursive map
record. Fetching one parent block does not authorize independent whole-block
writes from its children. Commit a single combined replacement.
\item Insert refreshed records/dummies into the pool and execute the fixed
eviction schedule: $2B$ lower subtrees per tree. Candidate selection and
pool cleanup retain the template costs derived below.
\item Compare the full query key with both candidate keys and all eight
cuckoo-stash keys. Return a shared value and shared hit bit.
\end{enumerate}

Fresh leaf positions are generated as $\rho_G\xor\rho_E$ inside GC, from
independent private random contributions. Their XOR costs no AND gates;
input labels and OTs are charged. Fresh positions are never opened when
assigned. The old position is disclosed only when that record is accessed
later. Untouched state labels persist; touched state obtains fresh output
labels. Tables are generated for the current work and streamed, not retained
as a complete multi-query garbled program.

Intersecting fetch paths require aggregation before clearing physical
slots. After deduplication, at most one representative targets any live
record. Per-slot deltas can therefore be XOR-aggregated and committed once.
Deterministically chosen eviction subtrees are disjoint below the pool;
this property does not make intersecting fetch writes independent.

\section{Deriving the circuit and communication ledger}
\label{sec:ledger}
We count ordinary Boolean ANDs separately from OT and shuffle bytes. With
FreeXOR and half gates, an AND contributes 32 bytes from $G$ to $E$, while an
XOR contributes no garbled table~\cite{halfgates}. FreeXOR still requires
local computation and memory traffic. This section makes explicit which
counts follow directly from kernels and which are composition allowances.

\subsection{Reusable Boolean kernels}
For a bitonic sort of $m$ items, let $K=2^a$ with
$a=\ceil{\log_2m}$. There are
\begin{equation}
 T(K)=Ka(a+1)/4
\end{equation}
compare-exchanges. For a $k$-bit sorting key, $p$ payload bits, and $v$ bits
of inverse routing through saved controls, the template charges
\begin{equation}
 A_{\rm sort}(m,k,p,v)=T(K)(4k-2+p+v).
\label{eq:sort}
\end{equation}
The comparison uses at most $3k-2$ ANDs; swapping the key costs $k$,
the payload costs $p$, and inverse routing costs $v$. Conditional swaps
already use one AND per bit: no additional factor-of-two discount applies.

Two work-efficient segmented sweeps, adjacent group comparisons, and final
masking are budgeted by
\begin{equation}
 A_{\rm group}(m,g,w)=(m-1)\max(g-1,0)
       +4(m-1)(w+2)+m(w+4).
\label{eq:group}
\end{equation}
The network and comparison formulas specify concrete kernels. Equation~
\eqref{eq:group}'s use in complete pool/duplicate phases is a composition
allowance; it is not a compiled gate trace. Define
\begin{equation}
 \Phi(m,k,g,w,v)=A_{\rm sort}(m,k,w+3,v)+A_{\rm group}(m,g,w).
\end{equation}

\subsection{Work retained at each recursive tree}
Suppress the level subscript and write $h,W,L$ for the current widths and
path length. The following counts are per table and per batch:
\begin{align}
 A_{\rm conflicts}&=A_{\rm sort}(B,h+\log_2B+1,W+2,W)
                       +A_{\rm group}(B,h,W),\\
 A_{\rm pathfetch}&=BZL(h+W+2),\\
 A_{\rm depth}&=2h(P+B),\\
 A_{\rm select}&=\Phi(P+3B,r+e+\ceil{\log_2(P+3B)}+2,r,W,W),\\
 A_{\rm cleanup}&=\Phi(P+3B,1+\ceil{\log_2(P+3B)},1,W,0),\\
 A_{\rm pathchoice}&=Bh,
\end{align}
where $e=\max(1,\ceil{\log_2(h+1)})$. Request-conflict work is scheduled
on both sides of retrieval: forward grouping first, result propagation and
inverse routing afterward. The old pool-retrieval sort is replaced in its
entirety by the shuffle protocol, rather than discounted a second time.

The lower-subtree eviction budget is
\begin{align}
 A_{\rm evict}={}&(2BZL)(2h)
  +2BL(Z-1)(3e-2+e+\ceil{\log_2Z}+1) \nonumber\\
 &+2B\cdot8L(6e+2)+2B(2Z+4)LW.
\label{eq:evict}
\end{align}
These terms are eligibility, bucket-winner selection, a parallel planner,
and record movement. The planner factor 8 and extra movement factor 4 are
exposed allowances inherited from the existing model. They must be checked
against an implemented parallel eviction circuit; Equation~\eqref{eq:evict}
is not asserted as an exact count from the Circuit OPRAM paper.

\subsection{Recursive updates, flat map, and application logic}
For a map block containing $F$ labels of $h$ bits, write $U=Fh$ and
$h'=h-\log_2F$. Coalescing updates is budgeted by
\begin{equation}
 A_{\rm packedmap}=B\bigl((F-2)+F+6U\bigr)
       +4(B-1)(U+F+2)+(B-1)\max(h'-1,0).
\end{equation}
Here ``packed'' means multiple scalar position labels in one map block;
it is unrelated to the two-label value encoding.

For the final flat map of $C$ cells with $h$-bit labels, the fused read/update
has cost
\begin{align}
 A_{\rm base}={}&A_{\rm sort}(B,\log_2C+\log_2B+1,h+2,h)
    +A_{\rm group}(B,\log_2C,h)\nonumber\\
 &+B\bigl((C-2)+C+Ch\bigr).
\end{align}
The last line pays for decoding, representative gating, and one shared
read/write product per label bit. It saves $BCh$ ANDs per table relative
to separate read and write masks. No unproved one-hot replacement is credited.

With two table candidates and eight cuckoo-stash entries, the application
lookup contributes
\begin{equation}
 A_{\rm app}=B\bigl(10(255+1+256)+9\bigr)=5129B.
\end{equation}
The 255 ANDs test a 256-bit equality, one gates occupancy, 256 mask the
returned value, and nine combine hit flags. A dictionary key is assumed
unique among live entries.

\subsection{The new pool term and directional traffic}
Let $t=\ceil{\log_2(B+1)}$ and
$b_\ell=\ceil{(W_\ell+t)/8}$. For each pool and batch, the model charges
\begin{align}
 A_{\rm poolmatch}&=BP(h_\ell+2),\\
 C_{\rm shuffle}&=2nb_\ell+4S(n)b_\ell+32S(n),\\
 C_{\rm tags}&=2\ceil{nt/8},\\
 I_{\rm reinput}&=nW_\ell\quad\text{shared bits}.
\end{align}
The $32S(n)$ term is an illustrative 16-byte OT-extension allowance for each
of the $2S(n)$ switches. Actual OT setup, padding, and implementation constants
must be measured. Shuffle corrections are paid for every fresh batch, even
when prepared offline. The tag vectors are bit-packed; rounding is per
message.

The primary model converts all shuffled records back to Boolean labels,
including the selected records and remaining pool. A shared-bit input uses
16 bytes for the garbler's share label and classic evaluator input OT
(32 forward, 16 reverse): 48 bytes forward and 16 reverse. This is a
conservative standard interface, not an optimal mixed-protocol conversion.
Tag bits are opened and do not need re-input. The shuffle and tag traffic is
symmetric across the two directions because the two private permutations
exchange the parties' roles.

Let $A$ be total ordinary ANDs per \emph{complete query}, including both
tables, application logic, and 10,240 CMAC ANDs; let $J$ be the corresponding
shared re-input bits, $U=2\sum_\ell h_\ell$ the fresh-position bits, $V=U$
the old-path bits, and $C_s$ the sum of shuffle and tag bytes per query.
Before framing,
\begin{align}
 C_{G\to E}&=32A+48J+\tfrac12C_s+48(256)+48U+V/8,\\
 C_{E\to G}&=16J+\tfrac12C_s+16(256)+16U+V/8.
\label{eq:wire}
\end{align}
Old paths are opened to both parties, costing $V/8$ bytes in each direction.
The query key is entered as two 256-bit shares. No supplied-index inputs are
charged, because indices are computed internally. No plaintext answer
decoding to $E$ is charged: the answer remains shared. Base OT, secret PRF
key setup, initialization, and authorized client delivery are separate
amortized or deployment-specific items. The reported wire counts multiply
each direction by 1.05 for framing.

\section{Numerical results and selection of parameters}
\label{sec:results}
\subsection{Batch and recursion search}
The search enumerates $B\in\{8,16,32,64,128,256\}$,
$F\in\{4,8,16,32\}$, and requested base cutoffs
$\{16,32,64,128,256,512,1024,2048,4096\}$: 216 combinations per database
size. Each row below selects the cheapest network/state/I/O configuration
at that batch size. Both final selections use $B=32$, $F=4$, and an actual
base map of 2048 entries. The requested cutoff may have several equivalent
values; the actual instantiated base is what matters. Small batches lose
amortization of the fixed 128-record pool slack. Larger batches grow the
remaining padded sorting networks. Consequently the objective is not
monotone in batch size.

\begin{center}\small
\begin{tabular}{lrrrrrr}\toprule
$B$ & $F_{28}/F_{32}$ & MB/q, 28 & MB/q, 32 & \$/yr, 28 & \$/yr, 32 & $\tau_B$ limit (ms)\\\midrule
8 & 4/4 & 342.7 & 429.9 & \$55.89M & \$71.97M & 0.56\\
16 & 4/4 & 395.4 & 494.2 & \$64.30M & \$82.26M & 1.12\\
32 & 4/4 & 237.8 & 301.3 & \$39.05M & \$51.40M & 2.24\\
64 & 4/8 & 273.3 & 342.2 & \$44.70M & \$57.83M & 4.48\\
128 & 8/8 & 314.9 & 392.5 & \$51.30M & \$65.88M & 8.96\\
256 & 8/8 & 366.0 & 454.3 & \$59.37M & \$75.60M & 17.92\\
\bottomrule\end{tabular}
\end{center}

Traffic includes both directions and 5\% framing;
prices exclude CPU sizing. The final column is the one-state batch reuse
limit, not a predicted service time. Padding effects and the fixed slack
make this an optimum of the selected implementation template, not a lower
bound for OPRAM.

\subsection{Instantiated recursive trees}
Every selected tree has a 317-record pool, $r=6$, and a shuffle of 349
records. The final map has 2048 scalar entries of 11 bits. The table lists
each tree's address width $h$, payload width $D$, complete record width $W$,
and lower-path length $L=h-5$.

\begin{center}\small
\begin{tabular}{lrrrrrrrr}\toprule
Level & $h_{28}$ & $D_{28}$ & $W_{28}$ & $L_{28}$ & $h_{32}$ & $D_{32}$ & $W_{32}$ & $L_{32}$\\\midrule
0 & 29 & 513 & 572 & 24 & 33 & 513 & 580 & 28\\
1 & 27 & 116 & 171 & 22 & 31 & 132 & 195 & 26\\
2 & 25 & 108 & 159 & 20 & 29 & 124 & 183 & 24\\
3 & 23 & 100 & 147 & 18 & 27 & 116 & 171 & 22\\
4 & 21 & 92 & 135 & 16 & 25 & 108 & 159 & 20\\
5 & 19 & 84 & 123 & 14 & 23 & 100 & 147 & 18\\
6 & 17 & 76 & 111 & 12 & 21 & 92 & 135 & 16\\
7 & 15 & 68 & 99 & 10 & 19 & 84 & 123 & 14\\
8 & 13 & 60 & 87 & 8 & 17 & 76 & 111 & 12\\
9 & 11 & 52 & 75 & 6 & 15 & 68 & 99 & 10\\
10 & -- & -- & -- & -- & 13 & 60 & 87 & 8\\
11 & -- & -- & -- & -- & 11 & 52 & 75 & 6\\
\bottomrule\end{tabular}
\end{center}

There are 10/12 dependent path-disclosure stages.
Choosing fanout 8 for the larger database offers a latency/storage tradeoff;
the selected fanout 4 minimizes the stated aggregate price objective, which
does not price round-trip latency.

\subsection{Ordinary garbled-circuit work}
The following rows are millions of ANDs per complete query, already summed
over both tables and every recursive level. They exclude OT and shuffle
messages, which are accounted separately.

\begin{center}\small
\begin{tabular}{lrr}\toprule
Component (million AND/query) & $2^{28}$ & $2^{32}$\\\midrule
Request grouping and private result routing & 0.082943 & 0.103014\\
Lower-path retrieval and clearing & 0.205776 & 0.288648\\
Pool UID matching and tags & 0.139480 & 0.182592\\
Eviction-candidate depth & 0.008725 & 0.011517\\
Pool eviction selection & 3.270576 & 4.008054\\
Pool cleanup & 1.723449 & 2.110749\\
Lower-subtree eviction & 1.479984 & 2.071104\\
Duplicate-masked path selection & 0.000400 & 0.000528\\
Coalesced recursive-map updates & 0.015771 & 0.021057\\
Fused flat position-map exchange & 0.054744 & 0.054744\\
Dictionary key match and value selection & 0.005129 & 0.005129\\
AES-CMAC query frontend & 0.010240 & 0.010240\\
\textbf{Total} & \textbf{6.997218} & \textbf{8.867375}\\
\bottomrule\end{tabular}
\end{center}

For $N=2^{28}$, eviction selection and pool cleanup account for 71.4\% of ordinary ANDs.

For $N=2^{32}$, eviction selection and pool cleanup account for 69.0\% of ordinary ANDs.

These retained phases dominate after retrieval has been
improved. Optimizing the shuffle alone cannot remove them.

\subsection{Pool protocol: all costs, not only online bytes}
The main pool terms below include both tables and all trees, per query.
Re-input is charged even for the data already moved by the shuffle; these
are distinct operations. The unframed total is comparable to the old pool
phase only, not to the entire query.

\begin{center}\small
\begin{tabular}{lrr}\toprule
Pool component (MB/query) & $2^{28}$ & $2^{32}$\\\midrule
Metadata match tables & 4.463360 & 5.842944\\
Online masked arrays & 0.009772 & 0.011997\\
Offline switch corrections & 0.147280 & 0.180812\\
OT-extension allowance & 0.052600 & 0.063120\\
Tag-share opening & 0.000327 & 0.000393\\
GC re-input, both directions & 2.343884 & 2.882740\\
Total pool retrieval & 7.017224 & 8.982006\\
\bottomrule\end{tabular}
\end{center}

\subsection{Complete directional communication}

\begin{center}\small
\begin{tabular}{lrr}\toprule
Traffic component (MB/query) & $2^{28}$ & $2^{32}$\\\midrule
Garbled tables, forward & 223.910968 & 283.756016\\
Pool re-input, forward & 1.757913 & 2.162055\\
Pool re-input, reverse & 0.585971 & 0.720685\\
Shuffle + tags, both directions & 0.209980 & 0.256322\\
Query shares, both directions & 0.016384 & 0.016384\\
Fresh-position inputs, both directions & 0.025600 & 0.033792\\
Old-path disclosure, both directions & 0.000100 & 0.000132\\
\midrule Forward, after framing & 237.095679 & 300.388127\\
Reverse, after framing & 0.736582 & 0.904529\\
\textbf{Total, after framing} & \textbf{237.832261} & \textbf{301.292656}\\
\bottomrule\end{tabular}
\end{center}

PRF setup is 6400 ANDs plus 8192 input bytes per
initialized state, before framing, and is not amortized into these per-query
rows. At an epoch of $T$ queries, add $(32\cdot6400+8192)/T$ bytes per query
before framing, together with any state handoff. Base OT and initial database
construction remain separate.

\subsection{Packed-value sensitivity without an unsupported gate discount}
A 256-bit ordinary Boolean value uses 4096 label bytes per party. A
two-carrier representation would use 32 bytes. For $M_{\rm value}$ physical
data/stash slots, the conditional state saving is exactly
$4064M_{\rm value}$ bytes per party. Recursive leaf labels are unchanged.

\begin{center}\small
\begin{tabular}{lrr}\toprule
State per party (TB) & $2^{28}$ & $2^{32}$\\\midrule
Primary Boolean representation & 64.699 & 1,061.578\\
Conditional two-carrier values & 38.517 & 642.665\\
\bottomrule\end{tabular}
\end{center}

This is a storage sensitivity, not a complete second
deployment estimate. A packed backend must specify how shared values enter
and leave its encoding and what a secret-controlled move costs. The present
shuffle transports ordinary XOR shares; the Boolean color-bit conversion does
not prove a conversion for two opaque carriers. The earlier hypothetical
``two half-gate equivalents per value move'' cannot be included without those
interfaces. Ordinary compact XOR-share storage is another valid design
direction, but needs re-input accounting on every subsequently touched field.


\section{Persistent state, I/O, and annual capacity cost}
\label{sec:deployment}
The number of persistent slots in a tree is
\begin{equation}
 Z\bigl(2^{h_\ell+1}-2^r\bigr)+P.
\end{equation}
Therefore the primary model's persistent logical bits are
\begin{equation}
 S_{\rm bits}=2\left[\sum_\ell
 \left\{Z(2^{h_\ell+1}-2^r)+P\right\}W_\ell+C h_{\rm last}\right]
 +8\cdot513+1536.
\end{equation}
Each party retains one 16-byte label per logical bit, so the total across
both parties is $32S_{\rm bits}$ bytes. The garbler needs one zero label plus
a global FreeXOR offset, not two independently stored labels per bit.
The final 1,536 bits retain the expanded AES round keys and CMAC subkey.
The newly added dummy rows and shuffle masks are batch workspace, not
persistent additional tree capacity.

The persistent-I/O envelope touches the whole pool, $B$ fetch paths, and
$2B$ eviction paths at every tree. It counts overlaps repeatedly and assumes
scratch intermediates remain in RAM:
\begin{equation}
 T_{\rm bits}=2\left[\sum_\ell(P+3BZL_\ell)W_\ell+C h_{\rm last}\right]
            +8\cdot513.
\end{equation}
Reading and writing label state costs $32T_{\rm bits}/B$ bytes per query
\emph{per party}. Converting pool labels to XOR shares does not imply that
the rest of the database has been converted to compact share storage.

\subsection{Historical deployment assumptions}
The dollar estimates reuse the repository's official AWS US East
(N.~Virginia) price snapshot, published 12 October 2026. They are not a new
price quote. The assumed c6in.32xlarge rate is \$7.2576/hour; each node has
256 GiB RAM, 200 Gbit/s network, and up to 100 Gbit/s EBS bandwidth.
The gp3 capacity rate is \$0.96/GiB-year~\cite{awsnetwork,awsebs,awsprice}.

Applying network-efficiency, distribution, and utilization factors gives
\begin{equation}
 c_{\rm NIC}=\frac{200\cdot10^9}{8}(0.65)(0.70)(0.70)
            =7.9625\cdot10^9\ \text{bytes/s per direction}.
\end{equation}
Full-duplex NIC capacity is enforced at both roles:
\begin{equation}
 n_{\rm role}=\ceil{\frac{Q\max(C_{G\to E},C_{E\to G})}{c_{\rm NIC}}},
 \qquad D_{\rm EC2}=2n_{\rm role}(8760)(7.2576).
\end{equation}
Both servers' worker fleets are in one availability zone; the model assigns
no marginal cross-AZ transfer charge in this placement. Intra-party data
distribution is not free in an implementation and is not explicitly routed
by this aggregate NIC calculation.

For storage, annual capacity cost is
$0.96\,S_{\rm bytes}/2^{30}$. The inherited I/O scenario uses 64 KiB aggregated
operations, provisioned with $0.65\cdot0.70$ throughput headroom, gp3 limits
of 2,000 MiB/s and 80,000 IOPS per volume, and at least four IOPS per provisioned
MiB/s. It prices extra throughput at \$0.48/(MiB/s)-year and extra IOPS at
\$0.06/IOPS-year after each volume's included 125 MiB/s and 3,000 IOPS.
This is a specified provisioning scenario, not a minimal price or a disk
latency guarantee. The model also checks volume attachments and EBS host
capacity, taking a maximum with the NIC fleet rather than double-counting
the same host.

\subsection{Selected annual budgets}
\begin{center}\small
\begin{tabular}{lrr}\toprule
Both-party cost unless stated & $2^{28}$ & $2^{32}$\\\midrule
NIC hosts per party & 298 & 378\\
Additional capacity/I/O EC2/year & \$0.00M & \$0.00M\\
Network-sized EC2/year & \$37.89M & \$48.06M\\
SSD capacity/year & \$0.12M & \$1.90M\\
Modeled persistent I/O/year & \$1.04M & \$1.43M\\
\textbf{Combined annual capacity budget} & \textbf{\$39.05M} & \textbf{\$51.40M}\\
Persistent R+W/query/party (MB) & 18.777 & 25.919\\
State capacity, both parties (TB) & 129.399 & 2,123.157\\
\bottomrule\end{tabular}
\end{center}

The minimum streaming buffer allowance is 64 MiB per
party. It is not the total RAM requirement: workers, in-flight circuit labels,
OT matrices, shuffle correlations, network buffers, and eviction scratch must
be provisioned. No full-horizon garbled program is retained.

\subsection{Known-work CPU sensitivity}
At the following assumed effective ordinary-AND rates per node, take the
maximum of NIC, storage/I/O, and ordinary-gate host counts. These figures
include the stated state and I/O budgets, but \emph{exclude} the additional
CPU cost of OT, shuffle, free-XOR, serialization, and distributed execution.
They therefore are not complete CPU-priced service budgets.

\begin{center}\small
\begin{tabular}{lrrrr}\toprule
Million AND/s/node & Hosts/party, 28 & Hosts/party, 32 & \$/yr, 28 & \$/yr, 32\\\midrule
50 & 1400 & 1774 & \$179.17M & \$228.90M\\
200 & 350 & 444 & \$45.66M & \$59.79M\\
1,000 & 298 & 378 & \$39.05M & \$51.40M\\
\bottomrule\end{tabular}
\end{center}


\section{Throughput, interaction, and unresolved costs}
\label{sec:limits}
Ordinary Yao evaluation needs no network round trip per AND layer. However,
old-path discovery, the two online private permutations, opening tags, and
re-entering shared state create real dependencies. The number of recursive
trees is only the count of dependent \emph{path-disclosure} barriers, not
the protocol's total interaction rounds. The two tables can proceed in
parallel; a concrete implementation must schedule the additional interfaces
and account for their RTTs.

Let $\tau_B$ be the measured time until one independently evolving ORAM
state can service its next batch, including preprocessing availability,
evaluation, updates, and synchronization. With $K$ independent states and
target utilization $u=0.70$, a necessary capacity condition is
\begin{equation}
 Q\le\frac{uKB}{\tau_B},\qquad
 K\ge\ceil{\frac{Q\tau_B}{uB}}.
\end{equation}
At $B=32$ and $Q=10{,}000$, one state requires $\tau_B\le2.24$ ms.
Mean batch-fill delay is approximately $(B-1)/(2Q)=1.55$ ms under regular
aggregate arrivals. Neither bound is established by buying enough network
nodes. A large distributed fleet does not automatically parallelize a
single dependent state transition.

Independent replicas multiply state capacity and initialization work.
They need independent hidden ORAM layouts; replicating the underlying
read-only dictionary does not authorize reusing correlated randomness or
garbled tables. Application writes would additionally require consistency
across replicas. The one-state annual figures must be revised if measured
$\tau_B$ requires additional states.

\paragraph{CPU accounting.}
An ordinary-AND rate can provide a workload sensitivity, but it does not
account automatically for OT extension, Waksman routing, PRG expansion,
label-to-share conversions, FreeXOR memory traffic, or distributed workers.
The report's rate sensitivities price known ordinary-gate work only; the
extra kernels still require benchmarking. Offline communication and CPU
must sustain the same long-run query rate as online evaluation.

\paragraph{What remains to be established.}
The match--tag--shuffle argument is conditional on the underlying OT,
garbling, and shuffle security, and on the unique-live-copy and duplicate
handling invariants. It is not a stand-alone proof for an implemented
Circuit OPRAM. Required next validations are: a full recursive-state
implementation; overflow and cuckoo-stash bounds at the desired lifetime;
an exact parallel eviction and packed-map-update circuit; hidden initial
placement; fresh-correlation and label-lifetime management; end-to-end
batch service time; and failure recovery. These are material limits of the
estimate, not priced deployment features.

\section{Reproduction and implementation checkpoints}
The accompanying Python model adds a new calculation rather than changing
the historical comparison files. From the repository root:
\begin{quote}\small\ttfamily
python -m mpc\_cost.interactive\_shuffle\_cost\\
python -m unittest mpc\_cost.test\_interactive\_shuffle -v\\
python -m mpc\_cost.render\_interactive\_shuffle\_report
\end{quote}
The JSON ledger is \path{mpc_cost/results/interactive_shuffle_comparison.json}.
The source model is \path{mpc_cost/interactive_shuffle_cost.py}; its inherited
kernel formulas are in \path{mpc_cost/two_party_batch.py}. The historical
price inputs are in \path{mpc_cost/data/aws_garbled_prices_2026-10-13.json}.
The PDF can be rebuilt from the self-contained delivered \LaTeX{} source by
running \texttt{pdflatex} twice. The source requires only standard packages
listed in its preamble and has an embedded bibliography.

Functional checks should verify that tags always form the same fixed
multiset, removal leaves exactly $P$ rows, dummies hide misses and duplicates,
and clear/reinsert operations preserve one live record per UID. Arithmetic
checks should independently reconcile each direction of traffic, all pool
terms, the flat-map saving, the PRF count, storage, and NIC rounding. These
checks validate the accounting and small functional kernels; they are not
cryptographic benchmarks or a security proof.

\begin{thebibliography}{99}\small
\bibitem{chan2017} T.-H. Hubert Chan and Elaine Shi.
\emph{Circuit OPRAM: Unifying Statistically and Computationally Secure ORAMs
and OPRAMs}. IACR ePrint 2016/1084, revised 2020.
\url{https://eprint.iacr.org/2016/1084}.
\bibitem{cgp2020} Melissa Chase, Esha Ghosh, and Oxana Poburinnaya.
\emph{Secret-Shared Shuffle}. ASIACRYPT 2020, ePrint 2019/1340.
\url{https://eprint.iacr.org/2019/1340}.
\bibitem{osn2021} Gayathri Garimella, Payman Mohassel, Mike Rosulek,
Saeed Sadeghian, and Jaspal Singh. \emph{Private Set Operations from
Oblivious Switching}. PKC 2021, ePrint 2021/243, Appendix A.3.
\url{https://eprint.iacr.org/2021/243}.
\bibitem{cop2025} Feng Han et al. \emph{Concretely Efficient Correlated
Oblivious Permutation}. IACR ePrint 2025/449.
\url{https://eprint.iacr.org/2025/449}.
\bibitem{halfgates} Samee Zahur, Mike Rosulek, and David Evans.
\emph{Two Halves Make a Whole: Reducing Data Transfer in Garbled Circuits
Using Half Gates}. EUROCRYPT 2015.
\url{https://www.cs.virginia.edu/~evans/pubs/ec2015/}.
\bibitem{heath2024} David Heath. \emph{Efficient Arithmetic in Garbled
Circuits}. EUROCRYPT 2024, ePrint 2024/139, Sections 4--5.
\url{https://eprint.iacr.org/2024/139}.
\bibitem{heath2021} David Heath and Vladimir Kolesnikov.
\emph{One Hot Garbling}. CCS 2021, ePrint 2022/798, Sections 4--5.
\url{https://eprint.iacr.org/2022/798}.
\bibitem{aesfips} NIST. \emph{Advanced Encryption Standard (AES)}.
FIPS 197, updated 2023. \url{https://doi.org/10.6028/NIST.FIPS.197-upd1}.
\bibitem{cmac} NIST. \emph{Recommendation for Block Cipher Modes of
Operation: The CMAC Mode for Authentication}. SP 800-38B, updated 2016.
\url{https://doi.org/10.6028/NIST.SP.800-38B}.
\bibitem{nistcircuits} NIST. \emph{Circuit Complexity: List of Circuits}.
Accessed 16 October 2026. The primary budget uses the classic 32-AND AES
S-box, not the newer 28-AND variants.
\url{https://csrc.nist.gov/projects/circuit-complexity/list-of-circuits}.
\bibitem{awsnetwork} AWS. \emph{Amazon EC2 C6i/C6in Instances};
\emph{EC2 Instance Network Bandwidth}.
\url{https://aws.amazon.com/ec2/instance-types/c6i/};
\url{https://docs.aws.amazon.com/AWSEC2/latest/UserGuide/ec2-instance-network-bandwidth.html}.
\bibitem{awsebs} AWS. \emph{Amazon EBS Pricing}; \emph{General Purpose SSD
Volumes}. \url{https://aws.amazon.com/ebs/pricing/};
\url{https://docs.aws.amazon.com/ebs/latest/userguide/general-purpose.html}.
\bibitem{awsprice} AWS. \emph{AmazonEC2 US East (N.~Virginia) Price List},
version 20261012193925, published 12 October 2026; locally retained snapshot.
\url{https://pricing.us-east-1.amazonaws.com/offers/v1.0/aws/AmazonEC2/20261012193925/us-east-1/index.csv}.
\end{thebibliography}
\end{document}
